\chapter*{Abstract}
\addcontentsline{toc}{chapter}{Abstract}

% ~300 words. Structure: Prior work → Gap → This thesis → Method → Results → Conclusion

This thesis extends a prior eight-month LoRaWAN indoor measurement campaign in which six
Arduino MKR WAN 1310 nodes collected 1,715,869 environmental observations across distinct
office locations. That campaign demonstrated that atmospheric pressure, CO\textsubscript{2}
concentration, temperature, humidity, and particulate matter (PM2.5) significantly influence
indoor LoRaWAN path loss, with a centralized XGBoost model achieving a coefficient of
determination $R^2 \approx 0.93$. However, the centralized design required continuous
transmission of raw sensor data, consumed excessive uplink bandwidth, offered no data
privacy, and produced a static model unable to adapt to post-deployment environmental
changes.

This thesis addresses these limitations by introducing a \gls{fl} framework combined with
\gls{tinyml}. Each MKR~WAN~1310 node runs a compact neural network
(Dense($5 \!\to\! 8 \!\to\! 1$), approximately 57~parameters,
$<$\,\SI{500}{\byte}) to predict expected path loss directly from local
environmental observations. Devices train the model on locally buffered samples and
communicate only quantized weight updates ($\approx$\,\SI{52}{\byte} per round) via
\gls{lorawan}. A central server aggregates updates using \gls{fedavg}, producing an
improved global model that is redistributed through rare downlinks.

Experimental evaluation on the real-world dataset compares centralized and federated
regression performance. The centralized baseline achieves $R^2 = $ [XX] with \gls{rmse}
$=$ \SI{[XX]}{\decibel}. The federated approach, simulated with non-\gls{iid} data
partitioned by physical device location (ED0--ED5), achieves $R^2 = $ [XX] after [N]
communication rounds---a reduction of only [X]\,\% relative to the centralized upper bound,
while reducing daily uplink volume by approximately $11\times$. Deployment on the
MKR~WAN~1310 confirms that the full inference pipeline fits within \SI{32}{\kilo\byte}
\gls{sram} and \SI{256}{\kilo\byte} Flash.

\noindent\textbf{Keywords:} Federated Learning, TinyML, LoRaWAN, Path Loss Modeling,
Edge Intelligence, Indoor Propagation, IoT

\clearpage

\chapter*{Zusammenfassung}
\addcontentsline{toc}{chapter}{Zusammenfassung}
[Deutsche Zusammenfassung hier einf\"ugen.]
\clearpage

\chapter*{Acknowledgments}
\addcontentsline{toc}{chapter}{Acknowledgments}
[Acknowledgments here.]
\clearpage
